% Exact proof for production Rethlas and source audit, 26 August 2026.

\begin{proof}
Put \(A=a+b\) and \(P=p+q\).  The delta-form module structure on
delta-currents is defined in \cite[\S6.4, equation~(6.4.1)]{GK17} by
\begin{equation}
 \langle\alpha\wedge T,\beta\rangle
 =(-1)^{AP}\langle T,\alpha\wedge\beta\rangle
 \label{eq:p21p3-module-pairing}
\end{equation}
for every compactly supported delta-form \(\beta\) of complementary
bidegree.  The current differentials are defined in \cite[\S6.9]{GK17} by
\begin{equation}
 \langle\mathrm d'T,\gamma\rangle
 =(-1)^{P+1}\langle T,\mathrm d'\gamma\rangle,
 \qquad
 \langle\mathrm d''T,\gamma\rangle
 =(-1)^{P+1}\langle T,\mathrm d''\gamma\rangle.
 \label{eq:p21p3-current-differential}
\end{equation}
The delta-forms themselves are a bigraded differential algebra by
\cite[\S4.13 and Proposition~4.15(iii)]{GK17}.

Let \(\beta\) have the bidegree complementary to
\(\mathrm d'(\alpha\wedge T)\).  From
\eqref{eq:p21p3-module-pairing} and
\eqref{eq:p21p3-current-differential},
\begin{equation}
 \langle\mathrm d'(\alpha\wedge T),\beta\rangle
 =(-1)^{A+P+1+AP}
   \langle T,\alpha\wedge\mathrm d'\beta\rangle.
 \label{eq:p21p3-left-sign}
\end{equation}
On the other hand,
\begin{align*}
 &\left\langle
   \mathrm d'\alpha\wedge T+(-1)^A\alpha\wedge\mathrm d'T,
   \beta\right\rangle\\
 &\quad=(-1)^{(A+1)P}
   \langle T,\mathrm d'\alpha\wedge\beta\rangle
   +(-1)^{A+A(P+1)+P+1}
   \langle T,\mathrm d'(\alpha\wedge\beta)\rangle.
\end{align*}
Expanding
\(
 \mathrm d'(\alpha\wedge\beta)
 =\mathrm d'\alpha\wedge\beta+(-1)^A
   \alpha\wedge\mathrm d'\beta
\), the two terms containing
\(\mathrm d'\alpha\wedge\beta\) cancel, while the remaining sign is
\((-1)^{A+P+1+AP}\).  Comparison with
\eqref{eq:p21p3-left-sign} proves
\eqref{eq:p21p3-dprime-module}.  Replacing \(\mathrm d'\) everywhere by
\(\mathrm d''\) gives \eqref{eq:p21p3-ddprime-module}; the total-degree
signs are identical.

Apply \eqref{eq:p21p3-dprime-module} to
\eqref{eq:p21p3-ddprime-module}.  Since \(\mathrm d''\alpha\) has total
degree \(A+1\), one obtains
\begin{align*}
 \mathrm d'\mathrm d''(\alpha\wedge T)
 ={}&\mathrm d'\mathrm d''\alpha\wedge T
   +(-1)^{A+1}\mathrm d''\alpha\wedge\mathrm d'T\\
 &+(-1)^A\mathrm d'\alpha\wedge\mathrm d''T
   +\alpha\wedge\mathrm d'\mathrm d''T,
\end{align*}
which is \eqref{eq:p21p3-dd-module}.  If
\(\mathrm d'T=\mathrm d''T=0\), then also
\(\mathrm d'\mathrm d''T=0\), and
\eqref{eq:p21p3-closed-module} follows.  Its form-level consequence is
immediate.

For the functorial premise, apply
\cite[Proposition~6.8]{GK17} to the generalized delta-form
\(\eta=\mathrm d'\mathrm d''\alpha\) and to zero.  That proposition says
precisely that equality of generalized delta-forms is detected by equality
of the associated ordinary currents after every algebraic pullback and on
every open subset lying over \(W\).  Hypothesis
\eqref{eq:p21p3-all-pullbacks} therefore gives
\(\mathrm d'\mathrm d''\alpha=0\) as a delta-form, so the already proved
case applies.

Finally, \cite[equation~(4.19.1)]{GK17} gives
\(B_{\mathrm{GK}}^{0,0}=A^{0,0}\).  Thus the module used above has only
smooth degree-zero multipliers.  Suppose now that \(a=b=0\) and
\eqref{eq:p21p3-smooth-single-current} holds.  Then
\(\mathrm d'\mathrm d''\alpha\) is a smooth form.  The local tropical
pairing sends it to the ordinary current in
\eqref{eq:p21p3-smooth-single-current}.  A boundaryless pure-dimensional
Berkovich space is very rich in the sense of Chambert--Loir--Ducros, and
the map from forms with tropically continuous coefficients to currents is
injective on such a space by \cite[Corollary~7.3.9]{CLD25}.  Applied on
the open set \(W\), this shows that
\eqref{eq:p21p3-smooth-single-current} implies
\(\mathrm d'\mathrm d''\alpha=0\) as a smooth form and hence as a
delta-form.  The closed-current case already proved gives the asserted
smooth degree-zero conclusion.

For general bidegree, Remark~6.7 and Proposition~6.8 of \cite{GK17} give
the all-pullback criterion used above, not a single-current detection
statement.  Piecewise-smooth functions instead lie in
\(P\!S^{0,0}=P^{0,0}\) by \cite[equation~(9.3.1)]{GK17}.  Section~9.5 of
\cite{GK17} constructs the algebra \(P\!S\!P\) and associates a
delta-current to each of its elements, but it does not extend the module
pairing of equation~(6.4.1) to a pairing of \(P\!S\!P\) with arbitrary
delta-currents.  In addition, the polyhedral differentials of an associated
piecewise-smooth delta-current have residue terms and are not compatible
with current differentiation in general \cite[Remark~9.6]{GK17}.  These
observations prove the scope assertions and keep the analytic question of
\cite[\S8.6.1]{CLD25} outside the conclusion.
\end{proof}
